% Options for packages loaded elsewhere
\PassOptionsToPackage{unicode}{hyperref}
\PassOptionsToPackage{hyphens}{url}
\documentclass[
]{article}
\usepackage{xcolor}
\usepackage[margin=1in]{geometry}
\usepackage{amsmath,amssymb}
\setcounter{secnumdepth}{-\maxdimen} % remove section numbering
\usepackage{iftex}
\ifPDFTeX
  \usepackage[T1]{fontenc}
  \usepackage[utf8]{inputenc}
  \usepackage{textcomp} % provide euro and other symbols
\else % if luatex or xetex
  \usepackage{unicode-math} % this also loads fontspec
  \defaultfontfeatures{Scale=MatchLowercase}
  \defaultfontfeatures[\rmfamily]{Ligatures=TeX,Scale=1}
\fi
\usepackage{lmodern}
\ifPDFTeX\else
  % xetex/luatex font selection
\fi
% Use upquote if available, for straight quotes in verbatim environments
\IfFileExists{upquote.sty}{\usepackage{upquote}}{}
\IfFileExists{microtype.sty}{% use microtype if available
  \usepackage[]{microtype}
  \UseMicrotypeSet[protrusion]{basicmath} % disable protrusion for tt fonts
}{}
\makeatletter
\@ifundefined{KOMAClassName}{% if non-KOMA class
  \IfFileExists{parskip.sty}{%
    \usepackage{parskip}
  }{% else
    \setlength{\parindent}{0pt}
    \setlength{\parskip}{6pt plus 2pt minus 1pt}}
}{% if KOMA class
  \KOMAoptions{parskip=half}}
\makeatother
\usepackage{color}
\usepackage{fancyvrb}
\newcommand{\VerbBar}{|}
\newcommand{\VERB}{\Verb[commandchars=\\\{\}]}
\DefineVerbatimEnvironment{Highlighting}{Verbatim}{commandchars=\\\{\}}
% Add ',fontsize=\small' for more characters per line
\usepackage{framed}
\definecolor{shadecolor}{RGB}{248,248,248}
\newenvironment{Shaded}{\begin{snugshade}}{\end{snugshade}}
\newcommand{\AlertTok}[1]{\textcolor[rgb]{0.94,0.16,0.16}{#1}}
\newcommand{\AnnotationTok}[1]{\textcolor[rgb]{0.56,0.35,0.01}{\textbf{\textit{#1}}}}
\newcommand{\AttributeTok}[1]{\textcolor[rgb]{0.13,0.29,0.53}{#1}}
\newcommand{\BaseNTok}[1]{\textcolor[rgb]{0.00,0.00,0.81}{#1}}
\newcommand{\BuiltInTok}[1]{#1}
\newcommand{\CharTok}[1]{\textcolor[rgb]{0.31,0.60,0.02}{#1}}
\newcommand{\CommentTok}[1]{\textcolor[rgb]{0.56,0.35,0.01}{\textit{#1}}}
\newcommand{\CommentVarTok}[1]{\textcolor[rgb]{0.56,0.35,0.01}{\textbf{\textit{#1}}}}
\newcommand{\ConstantTok}[1]{\textcolor[rgb]{0.56,0.35,0.01}{#1}}
\newcommand{\ControlFlowTok}[1]{\textcolor[rgb]{0.13,0.29,0.53}{\textbf{#1}}}
\newcommand{\DataTypeTok}[1]{\textcolor[rgb]{0.13,0.29,0.53}{#1}}
\newcommand{\DecValTok}[1]{\textcolor[rgb]{0.00,0.00,0.81}{#1}}
\newcommand{\DocumentationTok}[1]{\textcolor[rgb]{0.56,0.35,0.01}{\textbf{\textit{#1}}}}
\newcommand{\ErrorTok}[1]{\textcolor[rgb]{0.64,0.00,0.00}{\textbf{#1}}}
\newcommand{\ExtensionTok}[1]{#1}
\newcommand{\FloatTok}[1]{\textcolor[rgb]{0.00,0.00,0.81}{#1}}
\newcommand{\FunctionTok}[1]{\textcolor[rgb]{0.13,0.29,0.53}{\textbf{#1}}}
\newcommand{\ImportTok}[1]{#1}
\newcommand{\InformationTok}[1]{\textcolor[rgb]{0.56,0.35,0.01}{\textbf{\textit{#1}}}}
\newcommand{\KeywordTok}[1]{\textcolor[rgb]{0.13,0.29,0.53}{\textbf{#1}}}
\newcommand{\NormalTok}[1]{#1}
\newcommand{\OperatorTok}[1]{\textcolor[rgb]{0.81,0.36,0.00}{\textbf{#1}}}
\newcommand{\OtherTok}[1]{\textcolor[rgb]{0.56,0.35,0.01}{#1}}
\newcommand{\PreprocessorTok}[1]{\textcolor[rgb]{0.56,0.35,0.01}{\textit{#1}}}
\newcommand{\RegionMarkerTok}[1]{#1}
\newcommand{\SpecialCharTok}[1]{\textcolor[rgb]{0.81,0.36,0.00}{\textbf{#1}}}
\newcommand{\SpecialStringTok}[1]{\textcolor[rgb]{0.31,0.60,0.02}{#1}}
\newcommand{\StringTok}[1]{\textcolor[rgb]{0.31,0.60,0.02}{#1}}
\newcommand{\VariableTok}[1]{\textcolor[rgb]{0.00,0.00,0.00}{#1}}
\newcommand{\VerbatimStringTok}[1]{\textcolor[rgb]{0.31,0.60,0.02}{#1}}
\newcommand{\WarningTok}[1]{\textcolor[rgb]{0.56,0.35,0.01}{\textbf{\textit{#1}}}}
\usepackage{graphicx}
\makeatletter
\newsavebox\pandoc@box
\newcommand*\pandocbounded[1]{% scales image to fit in text height/width
  \sbox\pandoc@box{#1}%
  \Gscale@div\@tempa{\textheight}{\dimexpr\ht\pandoc@box+\dp\pandoc@box\relax}%
  \Gscale@div\@tempb{\linewidth}{\wd\pandoc@box}%
  \ifdim\@tempb\p@<\@tempa\p@\let\@tempa\@tempb\fi% select the smaller of both
  \ifdim\@tempa\p@<\p@\scalebox{\@tempa}{\usebox\pandoc@box}%
  \else\usebox{\pandoc@box}%
  \fi%
}
% Set default figure placement to htbp
\def\fps@figure{htbp}
\makeatother
\setlength{\emergencystretch}{3em} % prevent overfull lines
\providecommand{\tightlist}{%
  \setlength{\itemsep}{0pt}\setlength{\parskip}{0pt}}
\usepackage{bookmark}
\IfFileExists{xurl.sty}{\usepackage{xurl}}{} % add URL line breaks if available
\urlstyle{same}
\hypersetup{
  pdftitle={One-sample t-test and one-proportion test},
  hidelinks,
  pdfcreator={LaTeX via pandoc}}

\title{One-sample t-test and one-proportion test}
\author{}
\date{\vspace{-2.5em}}

\begin{document}
\maketitle

\textbf{Inference labs} use the five-step template: Hypothesis →
Visualise → Assumptions → Conduct → Conclude.

\subsection{Learning objectives}\label{learning-objectives}

\begin{itemize}
\tightlist
\item
  State a one-sample null and alternative hypothesis for a continuous
  outcome and for a proportion.
\item
  Run \texttt{t.test(mu\ =\ )} and \texttt{prop.test()} /
  \texttt{binom.test()} and interpret the output.
\item
  Apply the full five-step inference template in one working pass.
\end{itemize}

\subsection{Prerequisites}\label{prerequisites}

Labs 2.5, 3.3.

\subsection{Background}\label{background}

The one-sample \emph{t}-test asks whether the mean of a continuous
outcome differs from a pre-specified value μ₀. It is the inference
analogue of Table 1: a single-column calculation asking whether the
sample is consistent with a named population mean. The test statistic is
(x̄ − μ₀) / (s/√\emph{n}), compared to a \emph{t} distribution on
\emph{n} − 1 degrees of freedom.

The one-proportion test does the same for a binary outcome. When the
sample is large we use \texttt{prop.test()}, which is a Wald-type test
based on the normal approximation to the binomial; when the sample is
small, \texttt{binom.test()} gives an exact test based on the binomial
PMF. Both return a \emph{p}-value, a point estimate, and a CI.

These are the simplest tests in the course, but the five-step template
--- hypothesis, visualise, assumptions, conduct, conclude --- is the
same template used in every later inference lab. Getting the template
into muscle memory here pays dividends throughout the rest of the
course.

\subsection{Setup}\label{setup}

\begin{Shaded}
\begin{Highlighting}[]
\FunctionTok{library}\NormalTok{(tidyverse)}
\FunctionTok{set.seed}\NormalTok{(}\DecValTok{42}\NormalTok{)}
\FunctionTok{theme\_set}\NormalTok{(}\FunctionTok{theme\_minimal}\NormalTok{(}\AttributeTok{base\_size =} \DecValTok{12}\NormalTok{))}
\end{Highlighting}
\end{Shaded}

\subsection{1. Hypothesis}\label{hypothesis}

Scenario A (continuous): fasting blood glucose in a random sample of
\texttt{n\ =\ 40} patients from a clinic. Reference mean μ₀ = 95 mg/dL
in a healthy population. Two-sided test.

\begin{itemize}
\tightlist
\item
  H0: μ = 95
\item
  H1: μ ≠ 95
\item
  α = 0.05
\end{itemize}

Scenario B (proportion): of 120 patients offered a new screening
programme, \emph{k} = 78 accepted. Reference acceptance rate π₀ = 0.60.

\begin{itemize}
\tightlist
\item
  H0: π = 0.60
\item
  H1: π ≠ 0.60
\item
  α = 0.05
\end{itemize}

\subsection{2. Visualise}\label{visualise}

\begin{Shaded}
\begin{Highlighting}[]
\NormalTok{glucose }\OtherTok{\textless{}{-}} \FunctionTok{rnorm}\NormalTok{(n, }\AttributeTok{mean =} \DecValTok{100}\NormalTok{, }\AttributeTok{sd =} \DecValTok{12}\NormalTok{)}
\end{Highlighting}
\end{Shaded}

\begin{Shaded}
\begin{Highlighting}[]
\FunctionTok{tibble}\NormalTok{(glucose) }\SpecialCharTok{|\textgreater{}}
  \FunctionTok{ggplot}\NormalTok{(}\FunctionTok{aes}\NormalTok{(glucose)) }\SpecialCharTok{+}
  \FunctionTok{geom\_histogram}\NormalTok{(}\AttributeTok{bins =} \DecValTok{15}\NormalTok{, }\AttributeTok{fill =} \StringTok{"grey60"}\NormalTok{, }\AttributeTok{colour =} \StringTok{"white"}\NormalTok{) }\SpecialCharTok{+}
  \FunctionTok{geom\_vline}\NormalTok{(}\AttributeTok{xintercept =} \DecValTok{95}\NormalTok{, }\AttributeTok{linetype =} \DecValTok{2}\NormalTok{) }\SpecialCharTok{+}
  \FunctionTok{geom\_vline}\NormalTok{(}\AttributeTok{xintercept =} \FunctionTok{mean}\NormalTok{(glucose), }\AttributeTok{colour =} \StringTok{"firebrick"}\NormalTok{) }\SpecialCharTok{+}
  \FunctionTok{labs}\NormalTok{(}\AttributeTok{x =} \StringTok{"Fasting glucose (mg/dL)"}\NormalTok{, }\AttributeTok{y =} \StringTok{"Count"}\NormalTok{)}
\end{Highlighting}
\end{Shaded}

\begin{Shaded}
\begin{Highlighting}[]
\NormalTok{total  }\OtherTok{\textless{}{-}} \DecValTok{120}
\FunctionTok{tibble}\NormalTok{(}\AttributeTok{state =} \FunctionTok{c}\NormalTok{(}\StringTok{"accepted"}\NormalTok{, }\StringTok{"declined"}\NormalTok{),}
       \AttributeTok{count =} \FunctionTok{c}\NormalTok{(accept, total }\SpecialCharTok{{-}}\NormalTok{ accept))}
\end{Highlighting}
\end{Shaded}

\subsection{3. Assumptions}\label{assumptions}

\emph{t}-test: approximate normality of the sample mean, which --- by
the central limit theorem --- holds for \emph{n} = 40 unless the
underlying distribution is strongly skewed. Check with a Q-Q plot.

\begin{Shaded}
\begin{Highlighting}[]
\FunctionTok{tibble}\NormalTok{(glucose) }\SpecialCharTok{|\textgreater{}}
  \FunctionTok{ggplot}\NormalTok{(}\FunctionTok{aes}\NormalTok{(}\AttributeTok{sample =}\NormalTok{ glucose)) }\SpecialCharTok{+}
  \FunctionTok{stat\_qq}\NormalTok{() }\SpecialCharTok{+} \FunctionTok{stat\_qq\_line}\NormalTok{(}\AttributeTok{colour =} \StringTok{"firebrick"}\NormalTok{) }\SpecialCharTok{+}
  \FunctionTok{labs}\NormalTok{(}\AttributeTok{x =} \StringTok{"Theoretical"}\NormalTok{, }\AttributeTok{y =} \StringTok{"Sample"}\NormalTok{)}
\end{Highlighting}
\end{Shaded}

Proportion test: independent binary trials, same success probability.
For \texttt{prop.test} the large-sample normal approximation needs
\emph{np}₀ ≥ 5 and \emph{n}(1 − \emph{p}₀) ≥ 5; both are satisfied.

\subsection{4. Conduct}\label{conduct}

\begin{Shaded}
\begin{Highlighting}[]
\NormalTok{tt}
\end{Highlighting}
\end{Shaded}

\begin{Shaded}
\begin{Highlighting}[]
\NormalTok{pt\_exact }\OtherTok{\textless{}{-}} \FunctionTok{binom.test}\NormalTok{(}\AttributeTok{x =}\NormalTok{ accept, }\AttributeTok{n =}\NormalTok{ total, }\AttributeTok{p =} \FloatTok{0.60}\NormalTok{)}
\FunctionTok{list}\NormalTok{(}\AttributeTok{prop\_test =}\NormalTok{ pt\_large, }\AttributeTok{binom\_test =}\NormalTok{ pt\_exact)}
\end{Highlighting}
\end{Shaded}

\subsection{5. Concluding statement}\label{concluding-statement}

\begin{quote}
\textbf{Scenario A.} In a sample of \texttt{n} patients, mean fasting
glucose was \texttt{round(mean(glucose),\ 1)} mg/dL (SD
\texttt{round(sd(glucose),\ 1)}). The one-sample \emph{t}-test against
μ₀ = 95 gave \emph{t}(\texttt{tt\$parameter}) =
\texttt{round(tt\$statistic,\ 2)}, \emph{p} =
\texttt{signif(tt\$p.value,\ 2)}, mean difference
\texttt{round(mean(glucose)\ -\ 95,\ 1)} (95\% CI
\texttt{round(tt\$conf.int{[}1{]}\ -\ 95,\ 1)} to
\texttt{round(tt\$conf.int{[}2{]}\ -\ 95,\ 1)}).
\end{quote}

\begin{quote}
\textbf{Scenario B.} Acceptance was \texttt{accept} of \texttt{total}
(\texttt{round(100\ *\ accept\ /\ total,\ 1)}\%); the exact
one-proportion test against π₀ = 0.60 gave \emph{p} =
\texttt{signif(pt\_exact\$p.value,\ 2)}, 95\% CI
\texttt{round(pt\_exact\$conf.int{[}1{]},\ 2)} to
\texttt{round(pt\_exact\$conf.int{[}2{]},\ 2)}.
\end{quote}

The two scenarios share a template. The glossary changes --- mean and SD
for continuous, count and proportion for binary --- but the five steps
are the same.

\subsection{Common pitfalls}\label{common-pitfalls}

\begin{itemize}
\tightlist
\item
  Running a one-sided test to ``rescue'' a borderline two-sided
  \emph{p}-value.
\item
  Reporting the \emph{p}-value without the mean difference and its CI.
\item
  Using \texttt{prop.test} when the expected count of either outcome is
  very small; use \texttt{binom.test} there.
\end{itemize}

\subsection{Further reading}\label{further-reading}

\begin{itemize}
\tightlist
\item
  Rosner B. \emph{Fundamentals of Biostatistics}, chapter on one-sample
  inference.
\item
  Altman DG, \emph{Confidence intervals rather than p values}, BMJ 1986.
\end{itemize}

\subsection{Session info}\label{session-info}

\begin{Shaded}
\begin{Highlighting}[]

\end{Highlighting}
\end{Shaded}

\subsection{See also --- chapter index}\label{see-also-chapter-index}

\begin{itemize}
\tightlist
\item
  \href{../index.Rmd}{Methods in this chapter}
\item
  \href{../../../ch00-find-your-method.Rmd}{Find your method (Chapter
  0)}
\end{itemize}

\end{document}
